% Options for packages loaded elsewhere
\PassOptionsToPackage{unicode}{hyperref}
\PassOptionsToPackage{hyphens}{url}
\documentclass[
  brazilian,
  11pt,
  a4paper]{article}
\usepackage{xcolor}
\usepackage[margin=2.2cm]{geometry}
\usepackage{amsmath,amssymb}
\setcounter{secnumdepth}{-\maxdimen} % remove section numbering
\usepackage{iftex}
\ifPDFTeX
  \usepackage[T1]{fontenc}
  \usepackage[utf8]{inputenc}
  \usepackage{textcomp} % provide euro and other symbols
\else % if luatex or xetex
  \usepackage{unicode-math} % this also loads fontspec
  \defaultfontfeatures{Scale=MatchLowercase}
  \defaultfontfeatures[\rmfamily]{Ligatures=TeX,Scale=1}
\fi
\usepackage{lmodern}
\ifPDFTeX\else
  % xetex/luatex font selection
\fi
% Use upquote if available, for straight quotes in verbatim environments
\IfFileExists{upquote.sty}{\usepackage{upquote}}{}
\IfFileExists{microtype.sty}{% use microtype if available
  \usepackage[]{microtype}
  \UseMicrotypeSet[protrusion]{basicmath} % disable protrusion for tt fonts
}{}
\makeatletter
\@ifundefined{KOMAClassName}{% if non-KOMA class
  \IfFileExists{parskip.sty}{%
    \usepackage{parskip}
  }{% else
    \setlength{\parindent}{0pt}
    \setlength{\parskip}{6pt plus 2pt minus 1pt}}
}{% if KOMA class
  \KOMAoptions{parskip=half}}
\makeatother
\usepackage{graphicx}
\makeatletter
\newsavebox\pandoc@box
\newcommand*\pandocbounded[1]{% scales image to fit in text height/width
  \sbox\pandoc@box{#1}%
  \Gscale@div\@tempa{\textheight}{\dimexpr\ht\pandoc@box+\dp\pandoc@box\relax}%
  \Gscale@div\@tempb{\linewidth}{\wd\pandoc@box}%
  \ifdim\@tempb\p@<\@tempa\p@\let\@tempa\@tempb\fi% select the smaller of both
  \ifdim\@tempa\p@<\p@\scalebox{\@tempa}{\usebox\pandoc@box}%
  \else\usebox{\pandoc@box}%
  \fi%
}
% Set default figure placement to htbp
\def\fps@figure{htbp}
\makeatother
\ifLuaTeX
\usepackage[bidi=basic,shorthands=off,]{babel}
\else
\usepackage[bidi=default,shorthands=off,]{babel}
\fi
\ifLuaTeX
  \usepackage{selnolig} % disable illegal ligatures
\fi
\setlength{\emergencystretch}{3em} % prevent overfull lines
\providecommand{\tightlist}{%
  \setlength{\itemsep}{0pt}\setlength{\parskip}{0pt}}
\usepackage{amsmath,amssymb,booktabs}
\usepackage{fancyhdr}
\pagestyle{fancy}
\fancyhf{}
\fancyhead[L]{\small FLS 6183 | Lista 1}
\fancyhead[R]{\small Aulas 1 a 3}
\fancyfoot[C]{\thepage}
\setlength{\headheight}{14pt}
\setlength{\parskip}{5pt}
\renewcommand{\tablename}{Tabela}
\renewcommand{\figurename}{Figura}
\usepackage{bookmark}
\IfFileExists{xurl.sty}{\usepackage{xurl}}{} % add URL line breaks if available
\urlstyle{same}
\hypersetup{
  pdftitle={Lista 1: dos dados à previsão},
  pdfauthor={FLS 6183 - Métodos Quantitativos de Pesquisa II},
  pdflang={pt-BR},
  hidelinks,
  pdfcreator={LaTeX via pandoc}}

\title{Lista 1: dos dados à previsão}
\usepackage{etoolbox}
\makeatletter
\providecommand{\subtitle}[1]{% add subtitle to \maketitle
  \apptocmd{\@title}{\par {\large #1 \par}}{}{}
}
\makeatother
\subtitle{Probabilidade, distribuições, esperança condicional e BLP}
\author{FLS 6183 - Métodos Quantitativos de Pesquisa II}
\date{Conteúdo até a Aula 3 \textbar{} Revisão de 5 de setembro de 2026}

\begin{document}
\maketitle

\section{O que você vai investigar}\label{o-que-vocuxea-vai-investigar}

\textbf{Um resumo descreve bem os dados? Mais informação melhora a
previsão?} Você começará com a população dos municípios brasileiros e,
depois, retomará as votações de Brasil e China na Assembleia Geral das
Nações Unidas (AGNU). O percurso vai da importação e da descrição à
comparação de previsores.

Ao terminar, você deverá conseguir validar uma base, interpretar
momentos, distinguir probabilidades marginais e condicionais, simular
uma variável binária e comparar a função de esperança condicional (CEF)
com o melhor preditor linear (BLP).

\subsection{Organização e entrega}\label{organizauxe7uxe3o-e-entrega}

\begin{itemize}
\tightlist
\item
  Resolva os \textbf{exercícios 1 a 9}.
\item
  Trabalhe em três sessões: importação e descrição (1 a 4),
  probabilidade e previsão (5 e 6), BLP, exemplo não linear e síntese
  aplicada (7 a 9). As estimativas de tempo em cada bloco servem apenas
  para organizar o estudo.
\item
  Entregue pelo Moodle o PDF, a fonte \texttt{.Rmd} ou \texttt{.qmd} e
  um script \texttt{.R} com os cálculos. O documento deve carregar os
  resultados do script. Informe os integrantes conforme a modalidade de
  entrega anunciada no Moodle.
\item
  Apresente código, resultados selecionados e interpretação. Evite
  imprimir bases completas. Numere todas as figuras e tabelas e inclua
  legenda, unidade de medida e fonte.
\item
  Execute o trabalho em uma sessão nova do R antes de entregar. Use
  caminhos relativos e mantenha os arquivos de entrada intactos.
\item
  Consulte o Moodle para prazo e modalidade de entrega. Esta revisão não
  estabelece uma nova data.
\end{itemize}

\subsection{Arquivos e apoio}\label{arquivos-e-apoio}

Na pasta \texttt{dados/} do pacote estão dois arquivos:

\begin{itemize}
\tightlist
\item
  \texttt{populacao\_municipios\_2020.csv};
\item
  \texttt{brasil\_convergencia\_china\_1997\_2016.csv}.
\end{itemize}

O primeiro contém estimativas do IBGE para 1º de julho de 2020; o
segundo é o recorte didático usado nas Aulas 2 e 3. Use as cópias
fornecidas, sem baixar novamente as bases.

Tenha à mão os slides e os scripts das Aulas 2 e 3. Use explicitamente
\texttt{dplyr::select()} ao selecionar colunas. Caso recorra ao
\href{https://mgaldino.pythonanywhere.com/}{tutor}, registre a sua
tentativa, o diagnóstico e a verificação da solução. Você deve conseguir
explicar os comandos e as conclusões que entregar.

\textbf{Critério de qualidade:} cada resposta deve mostrar o que foi
calculado, em quais observações se baseia e o que o resultado permite
concluir.

\newpage

\section{1. Importar é também
verificar}\label{importar-uxe9-tambuxe9m-verificar}

\emph{Sessão 1, exercícios 1 a 4: aproximadamente 90 a 120 minutos.}

Importe \texttt{dados/populacao\_municipios\_2020.csv} para
\texttt{pop2020}.

\begin{enumerate}
\def\labelenumi{\alph{enumi}.}
\tightlist
\item
  Registre a primeira tentativa, mesmo se o comando rodar e produzir um
  resultado inesperado. Examine dimensões, nomes, tipos e algumas
  linhas. Se houver problema, formule uma hipótese, teste-a e explique a
  correção.
\item
  Valide os seguintes critérios: 5.570 linhas, três variáveis
  (\texttt{uf}, \texttt{nome\_munic}, \texttt{populacao}), 27 UFs,
  ausência de valores faltantes, população numérica e positiva e
  unicidade do par UF-nome do município. Verifique também a preservação
  de acentos e apóstrofos e eventuais problemas registrados pelo
  importador.
\item
  Explique a unidade de análise e a data de referência. Por que não
  basta verificar se há 5.570 linhas? Por que nomes de municípios,
  sozinhos, não são uma chave segura para a base nacional?
\item
  Crie uma cópia de trabalho: selecione as três colunas com
  \texttt{dplyr::select()}, renomeie \texttt{nome\_munic} para
  \texttt{municipio} e use \texttt{tolower()} para converter seus nomes
  em minúsculas. Preserve o objeto importado.
\end{enumerate}

\textbf{Produto:} diagnóstico breve e uma tabela numerada de
verificações, com critério, resultado observado e conclusão.

\section{2. Qual município estamos
descrevendo?}\label{qual-municuxedpio-estamos-descrevendo}

Crie um objeto com os municípios de São Paulo, mantendo a base nacional.
Neste exercício, defina \(Z\) como a população de um município paulista
escolhido com chances iguais. Se houver \(n\) municípios, chame de
\(z_i\) a população do município na linha \(i\) e de \(\bar z\) a média
desses valores. No R, você pode usar
\texttt{z\ \textless{}-\ sp\$populacao} para representar os valores
observados.

\begin{enumerate}
\def\labelenumi{\alph{enumi}.}
\tightlist
\item
  Quantos municípios há? Identifique o de menor população e os com mais
  de um milhão de habitantes. Apresente nomes e populações em uma
  tabela.
\item
  Calcule média, mediana, desvio-padrão e variância da população
  municipal. Informe a unidade de cada medida. Explique a diferença
  entre a população média de um município e a população total do estado.
\item
  Recalcule a média com a soma das populações dividida pelo número de
  municípios. A qual experimento aleatório essa média corresponde:
  escolher um município com chances iguais ou escolher um morador com
  chances iguais? Justifique sem efetuar uma nova ponderação.
\item
  Para os valores observados \(z_1,\ldots,z_n\) definidos acima, compare
  \(v_n=n^{-1}\sum_i(z_i-\bar z)^2\) com \texttt{var(z)}. Qual
  denominador o R usa? Verifique a relação entre os dois resultados. Use
  \(v_n\) quando tratar esta base como uma distribuição empírica com
  pesos iguais.
\end{enumerate}

\textbf{Pausa de interpretação:} uma estatística correta pode responder
à pergunta errada se a unidade de análise ou a ponderação estiverem
erradas.

\newpage

\section{3. Centro, dispersão e
forma}\label{centro-dispersuxe3o-e-forma}

Continue usando os municípios de São Paulo.

\begin{enumerate}
\def\labelenumi{\alph{enumi}.}
\tightlist
\item
  Produza uma figura com três painéis: distribuição da população de
  todos os municípios; distribuição dos municípios com menos de 50 mil
  habitantes; distribuição de \texttt{log(populacao)} para todos os
  municípios. Use histogramas ou densidades. Identifique as unidades e a
  população incluída em cada painel; não trate densidade como contagem.
\item
  Quantos municípios têm menos de 50 mil habitantes? Qual é sua
  proporção entre os municípios paulistas? Compare as distribuições e
  explique a diferença entre excluir municípios e transformar a escala.
\item
  Qual medida, média ou mediana, descreve melhor um município típico
  neste caso? Justifique com a figura e os números. A outra medida deixa
  de ser útil? Dê um exemplo de pergunta para a qual ela seria útil.
\item
  Defina \(W=Z/1000\), a população em milhares de habitantes. Antes de
  calcular, preveja como mudarão a média, o desvio-padrão e a variância.
  Verifique suas previsões no R e explicite as unidades resultantes.
\end{enumerate}

\textbf{Atenção à escala:} o logaritmo natural exige valores positivos e
muda as distâncias entre valores. Ele não altera quantos municípios
existem.

\section{4. Médias condicionais e média
nacional}\label{muxe9dias-condicionais-e-muxe9dia-nacional}

Retome todos os municípios brasileiros. Agora \(Z\) representa a
população de um município escolhido com chances iguais e \(U\)
representa sua UF.

\begin{enumerate}
\def\labelenumi{\alph{enumi}.}
\tightlist
\item
  Use \texttt{group\_by()} e \texttt{summarise()} para obter o número de
  municípios, a população total e a população média municipal em cada
  UF. Ordene a tabela pela média e identifique as UFs nos dois extremos.
\item
  Interprete \(E[Z\mid U=u]\). Uma UF com maior média municipal tem
  necessariamente maior população total? Sustente sua resposta com as
  colunas da tabela, sem confundir estado com município.
\item
  Reconstrua a média nacional usando as médias por UF ponderadas pela
  proporção de municípios de cada UF. Compare com a média nacional
  direta e com a média simples das 27 médias. Explique qual delas
  corresponde a escolher um município brasileiro com chances iguais.
\end{enumerate}

\textbf{Conexão com a Aula 3:} o item c é uma aplicação da Lei da
Esperança Iterada. As médias condicionais precisam ser combinadas com os
pesos da distribuição da variável em que condicionamos.

\newpage

\section{5. Dos votos às probabilidades e à
simulação}\label{dos-votos-uxe0s-probabilidades-e-uxe0-simulauxe7uxe3o}

\emph{Sessão 2, exercícios 5 e 6: aproximadamente 90 a 120 minutos.}

Importe a base da AGNU com \texttt{data.table::fread()}, como no
laboratório. Uma linha representa uma votação nominal com os votos de
Brasil e China observados. Defina \(Y=\texttt{convergente}\), que vale 1
se os votos coincidem e 0 caso contrário; e \(X=\texttt{pos\_2009}\),
que vale 0 em 1997--2008 e 1 em 2009--2016.

\textbf{Convenção para os exercícios 5 a 9:} trate as 1.762 linhas
válidas como uma distribuição empírica em que cada votação recebe o
mesmo peso, \(1/1.762\). Assim, probabilidades, esperanças e erros
quadráticos médios são calculados como médias sobre essas linhas. Não é
necessário sortear uma linha nem simular as votações reais; a simulação
pedida no item 5e é uma atividade separada. As medidas referem-se a
essas observações, não a uma população futura desconhecida.

\begin{enumerate}
\def\labelenumi{\alph{enumi}.}
\tightlist
\item
  Verifique a unicidade de \texttt{rcid}, ausentes nas variáveis usadas,
  anos entre 1997 e 2016, consistência entre \texttt{data} e
  \texttt{ano}, entre ano e período e valores 0/1 de \(X\) e \(Y\).
  Confira também se \(Y\) corresponde à igualdade dos dois votos. O
  ponto de conferência é \textbf{1.762 linhas válidas}. Não elimine
  linhas silenciosamente para atingir esse número. A coluna
  \texttt{voto\_importante} não será usada; seus valores ausentes não
  justificam excluir votações.
\item
  Construa uma tabela \(2\times2\) de contagens de \((X,Y)\) e, em
  seguida, a distribuição conjunta dividindo todas as células pelo total
  geral. Obtenha as duas distribuições marginais somando as células
  adequadas.
\item
  Calcule \(\Pr(Y=1\mid X=0)\), \(\Pr(Y=1\mid X=1)\) e
  \(\Pr(X=1\mid Y=1)\). Escreva em palavras a pergunta respondida por
  cada quantidade. Indique o denominador usado em cada caso.
\item
  Calcule \(p=E[Y]\), \(E[Y^2]\), \(\operatorname{V}(Y)\) e o
  desvio-padrão diretamente dos dados. Verifique \(E[Y^2]=E[Y]\) e
  \(\operatorname{V}(Y)=p(1-p)\). Use denominador \(n\) para a
  variância.
\item
  \textbf{Antes de simular}, preveja o centro de uma sequência de
  valores 0/1 com probabilidade \(p\) de valer 1. Fixe
  \texttt{set.seed(6183)} e gere 10.000 valores com
  \texttt{rbinom(10000,\ size\ =\ 1,\ prob\ =\ p)}. Compare a frequência
  de 1 e a variância empírica dos primeiros 100, primeiros 1.000 e todos
  os valores com \(p\) e \(p(1-p)\). Por que não precisam ser exatamente
  iguais?
\end{enumerate}

\textbf{Duas distinções:} aumentar a sequência não garante aproximação
monotônica em cada realização. A simulação cria sorteios independentes
por construção; isso não demonstra independência entre as votações
reais.

\newpage

\section{6. Qual informação melhora a
previsão?}\label{qual-informauxe7uxe3o-melhora-a-previsuxe3o}

Defina \(m(x)=E[Y\mid X=x]\). O número \(m(0)\) é a média de \(Y\) no
primeiro período; \(m(X)\) atribui a cada linha a média do seu período.
Para uma regra \(g\), o erro quadrático médio é \[
\operatorname{EQM}(g)=\frac{1}{n}\sum_{i=1}^{n}(Y_i-g(X_i))^2.
\]

\begin{enumerate}
\def\labelenumi{\alph{enumi}.}
\tightlist
\item
  Apresente uma tabela por período com número de votações, convergências
  e \(m(x)\). Calcule a diferença entre períodos em \textbf{pontos
  percentuais}. Explique a distinção entre a função \(m\), o número
  \(m(1)\) e a variável aleatória \(m(X)\).
\item
  Reconstrua \(E[Y]\) com \(m(0)\Pr(X=0)+m(1)\Pr(X=1)\). Por que a média
  simples das duas médias condicionais geralmente falha?
\item
  \textbf{Antes de executar}, ordene do menor para o maior EQM os
  previsores: constante 0,5; constante 1; média geral; média do período.
  Calcule os quatro EQMs nas mesmas 1.762 linhas e confira sua previsão.
  Use \texttt{left\_join()} para devolver a média de cada período às
  linhas; confira que a junção preservou o número de linhas e não gerou
  ausentes.
\item
  Defina \(e_i=Y_i-m(X_i)\). Verifique a média de \(e_i\) em cada
  período e no conjunto. Explique por que esses desvios têm média
  condicional zero.
\item
  Complete o argumento usando a Lei da Esperança Iterada (dica:
  condicione primeiro em \(X\); quais fatores passam a ser constantes?):
  \(E[(Y-m(X))(m(X)-g(X))]=0\). Explique por que isso implica \[
  \operatorname{EQM}(g)=\operatorname{EQM}(m)
  +E[(m(X)-g(X))^2].
  \] Em palavras, o que torna a CEF o melhor preditor sob perda
  quadrática?
\end{enumerate}

\textbf{Produto:} uma tabela de EQMs e um parágrafo distinguindo ser o
melhor preditor de produzir uma grande melhora. Quantifique a redução
absoluta do EQM ao passar da média geral à média por período.

\textbf{Limite da análise:} as médias foram construídas e avaliadas nos
mesmos dados. Um EQM menor aqui não garante melhor previsão de novas
votações e não identifica um efeito causal de 2009.

\newpage

\section{7. O que ganhamos e perdemos com uma
reta?}\label{o-que-ganhamos-e-perdemos-com-uma-reta}

\emph{Sessão 3, exercícios 7 a 9: aproximadamente 120 a 150 minutos.}

O melhor preditor linear (BLP) usa \(g(X)=\alpha+\beta X\), com \[
\beta=\frac{\operatorname{Cov}(X,Y)}{\operatorname{V}(X)},
\qquad \alpha=E[Y]-\beta E[X],
\] quando \(\operatorname{V}(X)>0\). \(\alpha\) é o intercepto e
\(\beta\) é a mudança na previsão quando \(X\) aumenta uma unidade.

\begin{enumerate}
\def\labelenumi{\alph{enumi}.}
\tightlist
\item
  Trabalhe com os valores observados \(x_i\) e \(y_i\) de \(X\) e \(Y\).
  Primeiro, escreva as fórmulas que \texttt{cov(x,\ y)} e
  \texttt{var(x)} usam, defina \(\bar x\) e \(\bar y\) e indique o
  denominador de cada fórmula. Depois, calcule a covariância e a
  correlação e explique, em palavras, o que cada medida resume. Por fim,
  explique por que podemos escrever o quociente
  \texttt{cov(x,\ y)\ /\ var(x)} sem mostrar o tamanho \(n\) da base:
  embora cada fórmula dependa de \(n\) e use o denominador \(n-1\), esse
  fator é o mesmo e se cancela no quociente acima.
\item
  Calcule \(\alpha\) e \(\beta\) pelas fórmulas, \textbf{sem usar
  \texttt{lm()}}. Obtenha \(g(0)\) e \(g(1)\) e compare com \(m(0)\) e
  \(m(1)\). Mostre que \(\beta=m(1)-m(0)\) e interprete a inclinação em
  pontos percentuais.
\item
  Verifique que a reta passa por \((E[X],E[Y])\). Explique por que, com
  \(X\) binária, dois coeficientes bastam para reproduzir os dois
  centros condicionais. Compare o EQM do BLP com o da CEF por período.
\end{enumerate}

\textbf{Pausa:} a seção seguinte usa uma distribuição conhecida para
separar covariância, independência, CEF e BLP. Não pede estimação de
MQO, erros-padrão, testes de hipótese ou interpretação causal.

\newpage

\section{8. Covariância zero não implica
independência}\label{covariuxe2ncia-zero-nuxe3o-implica-independuxeancia}

Considere o experimento aleatório em que a variável aleatória \(A\)
assume \(-1\), \(0\) e \(1\), cada um com probabilidade \(1/3\). Defina
a variável aleatória \(B=A^2\). Todas as probabilidades e esperanças
desta questão referem-se a essa distribuição exata; não é necessário
simular.

\begin{enumerate}
\def\labelenumi{\alph{enumi}.}
\tightlist
\item
  Faça uma tabela com os valores possíveis de \((A,B)\) e suas
  probabilidades. Calcule \(E[B\mid A=a]\) para cada \(a\in\{-1,0,1\}\),
  escreva a função \(m(a)=E[B\mid A=a]\) e calcule o EQM de \(m(A)\).
  Por que esse preditor acerta todos os valores neste exemplo?
\item
  Calcule \(E[A]\), \(E[B]\) e \(\operatorname{Cov}(A,B)\). Mostre cada
  etapa da soma e explique por que a covariância é zero.
\item
  Calcule o BLP \(g(A)=\alpha+\beta A\) usando
  \(\beta=\operatorname{Cov}(A,B)/\operatorname{V}(A)\) e
  \(\alpha=E[B]-\beta E[A]\). Encontre \(g(-1)\), \(g(0)\) e \(g(1)\);
  depois, calcule seu EQM e compare-o com o EQM da CEF. O que essa
  comparação ensina sobre uma reta quando a relação entre as variáveis
  não é linear?
\item
  \(A\) e \(B\) são independentes? Calcule \(\Pr(B=0\mid A=0)\) e
  \(\Pr(B=0)\) e use a comparação para justificar sua resposta. Não
  conclua independência apenas porque a covariância é zero.
\item
  Ainda no exemplo, defina \(u=B-(\alpha+\beta A)\). Compare \(E[u]\) e
  \(\operatorname{Cov}(A,u)\) com \(E[u\mid A=a]\). Explique por que as
  propriedades dos desvios do BLP são mais fracas que as da CEF.
\end{enumerate}

\textbf{Pausa:} covariância zero descreve ausência de associação linear,
não independência. A CEF pode capturar uma relação que o BLP não
captura.

\newpage

\section{9. Tema e período: mais detalhe, quais
limites?}\label{tema-e-peruxedodo-mais-detalhe-quais-limites}

Retome a AGNU. Agora \(T\) é a coluna \texttt{tema}. Trate cada texto
completo dessa coluna como uma categoria, inclusive combinações
separadas por ponto e vírgula. Assim, cada votação pertence a exatamente
uma categoria nesta análise. Mantenha a categoria sem codificação
temática.

\begin{enumerate}
\def\labelenumi{\alph{enumi}.}
\tightlist
\item
  Calcule a CEF empírica \(E[Y\mid T,X]\) e o número de votações para
  cada célula tema-período. Confira que a soma das contagens é 1.762.
  Identifique a menor célula e explique a fragilidade de uma média
  baseada em poucas observações.
\item
  Construa uma figura comparando as duas médias de cada tema, usando
  somente temas com pelo menos 30 votações \textbf{em cada um dos dois
  períodos}. Escolha pontos pareados ou diferenças em pontos
  percentuais. Informe o filtro na legenda e use eixos que deixem clara
  a comparação.
\item
  Devolva \textbf{todas} as médias tema-período às 1.762 linhas e
  calcule o EQM. O filtro do gráfico não deve restringir esta avaliação.
  Compare com os EQMs da média geral e da CEF por período. Explique por
  que subdividir grupos não aumenta o EQM quando as médias são
  calculadas e avaliadas nas mesmas linhas.
\item
  Se uma votação com dois temas passar a ocupar duas linhas, uma em cada
  tema, o que acontece com a unidade de análise e com a soma das
  contagens? Explique por que esse procedimento não pode ser usado
  silenciosamente na comparação de EQMs do item c.
\item
  Escolha um tema exibido e escreva \textbf{100 a 150 palavras}
  respondendo: como a convergência varia entre períodos? Inclua as duas
  médias, a diferença em pontos percentuais, os dois denominadores, o
  ganho preditivo geral de incluir tema e duas limitações: a avaliação
  nos mesmos dados e a ausência de identificação causal.
\end{enumerate}

\textbf{Antes de entregar:} releia o parágrafo sem olhar o código. Ele
informa qual comparação foi feita, em quais observações e com quais
limites?

\subsection{Conferência final do
trabalho}\label{conferuxeancia-final-do-trabalho}

\begin{itemize}
\tightlist
\item
  Os arquivos originais continuam intactos e os caminhos são relativos.
\item
  As verificações de dados aparecem antes dos resultados.
\item
  Todas as tabelas e figuras têm número, legenda e fonte.
\item
  Probabilidades, proporções e pontos percentuais estão distinguidos.
\item
  O script completo roda em uma sessão nova e gera os resultados do PDF.
\end{itemize}

\subsection{Materiais de referência}\label{materiais-de-referuxeancia}

Aronow e Miller (2019), \emph{Foundations of Agnostic Statistics},
capítulo 1 e seções 2.1--2.2.5; slides e laboratórios das Aulas 2 e 3.
Dados municipais: IBGE, estimativas para 1º de julho de 2020, arquivo
didático preservado. Votações: base didática AGNA, construída a partir
de \texttt{unvotes} 0.3.0, recorte Brasil-China, 1997--2016.

\end{document}
